\chapter{Open Questions}
\label{app:openq}

The framework is internally coherent but has many unresolved questions. This
appendix catalogs them, both for honesty about the framework's current state and
to guide future work.

\section{Foundational physics}

\subsection{Q1. Precise value of Cartasis density}
Current estimates place $\rhoc \sim 10^{50}\,\mathrm{kg\,m^{-3}}$, but this
varies by factors of 10--100 depending on coefficient choices in the
Einstein--Cartan Lagrangian. A clean derivation pinning the value would be
valuable.

\subsection{Q2. Is Cartasis density universal?}
Does the bounce happen at the same density for different matter content
(different fermion species, different bound states)? Or is the bounce condition
matter-dependent?

\subsection{Q3. How does Cartasis interact with rotation?}
Most analyses assume spherical (Schwarzschild) bounces. Real astrophysical black
holes are typically rotating (Kerr). The Kerr--Cartan bounce dynamics may differ
qualitatively from Schwarzschild--Cartan and might preferentially produce
specific chirality dynamics. Detailed analysis is needed.

\subsection{Q4. Quantum gravity at Cartasis}
The framework treats Einstein--Cartan classically. At
$\rhoc \sim 10^{50}\,\mathrm{kg\,m^{-3}}$, full quantum-gravity effects should be
subdominant (still 46 orders of magnitude below Planck density), but they might
matter for fine details. The full quantum-gravitational picture is unknown.

\subsection{Q5. Baryon-number conservation}
At trans-electroweak temperatures, sphaleron processes violate baryon number. At
higher temperatures (approaching Cartasis), more exotic violations might occur.
How conserved is $B$ across the bounce? This affects both the baryon-asymmetry
mechanism and the dark-matter interpretation.

\subsection{Q6. Beyond-Standard-Model physics}
The Standard Model is presumably an effective theory valid up to some scale.
Above that scale, new physics (SUSY? extra dimensions? unknown new fields?) might
dominate. Whether \SCT is robust to such new physics, or whether it requires the
SM to be the full story up to Cartasis, is unclear.

\section{Bounce dynamics}

\subsection{Q7. Flip vs.\ filter quantitative balance}
We have argued that the bounce involves both a CPT-like flip and a chirality
filter, with flip dominating. Quantitative work is needed: what is the ratio? How
does it depend on bounce parameters? Can it be measured?

\subsection{Q8. Bounce efficiency}
What fraction of the parent's mass--energy is transmitted to the baby's expanding
phase? Some might remain in an intermediate region, some might be lost to
gravitational waves, some might emerge as Hawking radiation in the parent's
exterior. The efficiency factor is currently a free parameter.

\subsection{Q9. Information flow through the bounce}
We have argued that information is preserved globally through entanglement
structures across the membrane. The specific quantitative details---what gets
entangled with what, how the Page curve is realized in bounce cosmology---are not
worked out.

\subsection{Q10. Persistent channel vs.\ severed connection}
We favor the position that the channel persists after the first bounce over the
position that each bounce severs it. But the detailed geometry of a persistent
channel that maintains across the parent black hole's lifetime needs rigorous
treatment.

\section{Supraverse structure}

\subsection{Q11. Equilibrium establishment}
We argue the supraverse is in thermodynamic equilibrium. How long does
establishment take? Is the supraverse currently in equilibrium or still
approaching it? What sets the equilibrium timescale?

\subsection{Q12. OG bounce rate}
What is the rate of OG bounce nucleation per unit supraverse spacetime? This
determines whether the supraverse contains many universes or few, and the
dependence on supraverse parameters.

\subsection{Q13. Universe size distribution shape}
We conjecture the distribution peaks near our observed size. What is the actual
shape? How wide is the distribution? Are most observers in universes very close
to our parameters, or is the distribution wide enough that observers could exist
in significantly different universes?

\subsection{Q14. Lineage interaction}
Do separate lineages (descended from different OG bounces) ever interact? If so,
what happens at boundaries? Can matter-lineage and antimatter-lineage universes
be near each other in the supraverse?

\subsection{Q15. Initial conditions for the supraverse}
The framework requires the supraverse to be old enough for equilibrium
establishment, but does not specify what the supraverse ``started from.'' Was
there an initial state, or is the supraverse genuinely eternal? This is partly a
question about the meaning of time at the supraverse level.

\subsection{Q15a. Do OG universes grow forever?}
Black holes have negative heat capacity, so above
$M_{\mathrm{crit}}=\hbar c^3/(8\pi G k_B T_{\mathrm{bath}})$ a universe-hole
accretes faster than it evaporates and grows. Whether an individual OG grows
without bound (runaway in a static infinite bath) or is regulated (by environment
depletion, by the void's expansion diluting the bath, by the stationary balance
of the foam) is open, and is entangled with the measure problem below.

\subsection{Q15b. The measure over infinitely many OG universes}
Axiom~\ref{ax:qm} plus infinity implies infinitely many OG nucleations. Making
``we are typical'' meaningful requires a regulated measure over that infinite,
possibly ever-growing ensemble---the measure problem shared with all multiverse
frameworks, here sharpened by ongoing growth (Q15a). Unsolved.

\subsection{Q15c. The cosmological constant in the void}
An infinite, eternally active quantum vacuum makes acute the question of why the
void is nearly flat rather than de Sitter at the QFT cutoff. The dark-energy
\emph{evolution} is reattributed to parent accretion (Chapter~\ref{chap:darksector}),
but the bare vacuum-energy problem is untouched, arguably worsened.

\section{Observational connection}

\subsection{Q16. CMB acoustic peaks}
$\Lambda$CDM predicts specific acoustic peaks in the CMB power spectrum that
match observations beautifully. Can \SCT reproduce these from bounce dynamics?
This is the central calculational test.

\subsection{Q17. Hubble tension resolution}
We claim \SCT might resolve the Hubble tension through modified early-universe
dynamics. The actual calculation showing this has not been done. If it cannot be
done, the framework needs revision.

\subsection{Q18. Dark-matter density profile}
Observed dark matter has specific clustering properties (NFW profiles,
mass--velocity relations). Can these be derived from parent gravitational
influence through Cartasis? Or do they imply something else?

\subsection{Q19. Galaxy-rotation asymmetry consistency}
If parent rotation imprints galaxy-rotation asymmetry, the relationship should be
specific and predictable. Shamir's observed asymmetry should match a particular
parent rotation axis. Whether this consistency holds at higher statistical
precision is unknown.

\subsection{Q20. Specific BBN constraints}
Big Bang Nucleosynthesis predicts light-element abundances. Does bounce cosmology
reproduce these? If the bounce dynamics differ from inflation $+$
radiation-domination, BBN constraints might require modifications.

\section{Theoretical issues}

\subsection{Q21. Connection to string theory / loop quantum gravity}
How does \SCT relate to other proposed extensions of GR? Could string theory or
LQG be ways of realizing the bounce in different formal frameworks? Are they
competitors or complements?

\subsection{Q22. Conformal cyclic cosmology relationship}
Penrose's CCC is also a non-singular cyclic cosmology~\citep{penrose2010}. How
does \SCT relate to CCC? Different but related, or distinct frameworks?

\subsection{Q23. ER\,$=$\,EPR specific realization}
The bounce structure suggests a geometric realization of ER\,$=$\,EPR. What is
the precise mathematical statement? Can it be cleanly derived?

\subsection{Q24. Holographic principle in the supraverse}
The holographic principle suggests that bulk physics can be encoded on
boundaries. Does the supraverse have a holographic structure? What is the
appropriate boundary? Is each Cartasis membrane a holographic surface for its
enclosed baby?

\subsection{Q25. Anthropic considerations}
We have implicitly used anthropic reasoning (observers exist in viable
universes). Can this be made more precise? Does the framework predict where
observers are in the supraverse population in a quantitatively testable way?

\section{Philosophical / interpretive}

\subsection{Q26. What is the supraverse ``embedded in''?}
Mathematically, the supraverse is a 4D manifold with non-trivial topology.
Physically, what is its ontological status? Is it the totality of existence, or
is there a meta-supraverse containing it?

\subsection{Q27. Connection between local quantum measurements and bounce
  dynamics}
We argued that gravitational decoherence is the small-scale version of bounce
dynamics. The precise relationship needs working out---is every quantum
measurement literally a miniature bounce-like splitting? Or just analogous?

\subsection{Q28. Implications for the meaning of physical constants}
If we are typical inhabitants of a supraverse population, our physical constants
are typical for our class of universe. Does this mean physical constants are
statistical features of the supraverse distribution, or are they fundamental? Or
is the distinction meaningless?

\subsection{Q29. Time at the supraverse level}
We argue the supraverse is static (a block universe). But the supraverse also
``supports'' universe nucleation and equilibrium establishment, which seem to
require time. How do we reconcile a static block universe with these
dynamical-sounding processes?

\section{Methodological}

\subsection{Q30. Falsifiability completeness}
The framework makes predictions (Chapter~\ref{chap:tests}). Are these
collectively sufficient to falsify the framework if it is wrong? Or could the
framework be twisted to accommodate any observational result? A genuinely
scientific framework should be unambiguously falsifiable. We have argued \SCT is,
but this should be checked carefully.

\subsection{Q31. Naturalness arguments}
We have argued \SCT is ``natural'' because it requires few postulates. But
naturalness arguments have led physics astray before (e.g.\ naturalness arguments
for SUSY have not worked out). Is \SCT's parsimony actually evidence for its
truth, or just for its appeal?

\subsection{Q32. Relationship to existing alternative cosmologies}
Many alternatives to $\Lambda$CDM exist: timescape (Wiltshire), inhomogeneous
cosmology, modified gravity (MOND/MOG), cyclic universes. \SCT shares features
with some of these. Detailed comparison would clarify what \SCT uniquely
contributes and what it borrows.

\section{Practical / programmatic}

\subsection{Q33. Publication strategy}
If/when this framework is written up rigorously, where would it best be
published? Specialized journals (\emph{Physical Review D}, \emph{Classical and
Quantum Gravity})? Or as a book/monograph?

\subsection{Q34. Collaboration approach}
The framework benefits from multiple expertises: numerical relativity, cosmology,
quantum field theory, philosophy of physics. How to assemble the right team for
serious development?

\subsection{Q35. Funding considerations}
This is theoretical/computational work that does not require large facilities.
Modest computational resources, person-time. What funding mechanisms are
appropriate?

\section{Priority ordering}

The most important open questions for near-term progress:

\begin{enumerate}
  \item \textbf{Q16 (CMB acoustic peaks):} decisive test; must be addressed for
    the framework to be taken seriously.
  \item \textbf{Q7 (flip vs.\ filter):} resolves chirality dynamics, enables the
    dark-matter prediction.
  \item \textbf{Q1 (precise $\rhoc$):} enables quantitative bounce calculations.
  \item \textbf{Q12--Q13 (supraverse statistics):} tests observer-placement
    predictions.
  \item \textbf{Q17 (Hubble tension):} high-impact if it works.
\end{enumerate}

Items 1--3 are calculable now with current Einstein--Cartan understanding plus
standard numerical methods. Items 4--5 require more theoretical development.
